\begin{bpnew}{thmborder}{Theorem: Uniqueness of Limits}
% ---------------------------------------------------------
% Theorem: Uniqueness of Limits
% Label: thm:uniqueness-of-limits
% ---------------------------------------------------------
\begin{theorem}[Uniqueness of Limits]
\label{thm:uniqueness-of-limits}

Let $(x_n) \subseteq \mathbb{R}$. If $x_n \to L$ and $x_n \to K$, then
$L = K$.

\smallskip
\noindent
\hyperref[prf:uniqueness-of-limits]{\textit{Go to proof.}}
\end{theorem}
\end{bpnew}

\begin{remark*}[Standard quantified statement]
\[
\begin{aligned}
&\forall (x_n) \subseteq \mathbb{R} \;\forall L \in \mathbb{R}
\;\forall K \in \mathbb{R}\\
&\quad\Bigl(
    \bigl(x_n \to L \;\land\; x_n \to K\bigr)
    \;\Longrightarrow\;
    L = K
\Bigr).
\end{aligned}
\]
\end{remark*}

\begin{remark*}[Theorem predicate reading]
\[
\operatorname{UniqueSequentialLimit}(x_n,L,K)
\;:=\;
\bigl(\operatorname{ConvergentSequence}(x_n,L)
\land
\operatorname{ConvergentSequence}(x_n,K)\bigr)
\Longrightarrow L=K.
\]
\end{remark*}

\begin{remark*}[Negated quantified statement]
\[
\begin{aligned}
&\exists (x_n) \subseteq \mathbb{R} \;\exists L \in \mathbb{R}
\;\exists K \in \mathbb{R}\\
&\quad\Bigl(
    \bigl(x_n \to L \;\land\; x_n \to K\bigr)
    \;\land\;
    L \neq K
\Bigr).
\end{aligned}
\]
\end{remark*}

\begin{remark*}[Negation predicate reading]
\[
\lnot\operatorname{UniqueSequentialLimit}(x_n,L,K)
\;:=\;
\operatorname{ConvergentSequence}(x_n,L)
\land
\operatorname{ConvergentSequence}(x_n,K)
\land L\neq K.
\]
\end{remark*}

\begin{remark*}[Failure modes]
The only failure mode is the existence of one real sequence with two
distinct real limits. In $\mathbb{R}$ this is impossible because
distinct real numbers can be separated by disjoint neighborhoods.
\end{remark*}

\begin{remark*}[Failure mode decomposition]
\[
\underbrace{\exists (x_n) \subseteq \mathbb{R}}_{\text{some sequence}}
\;\underbrace{x_n \to L \;\land\; x_n \to K}_{\text{converges to two values}}
\;\underbrace{L \neq K}_{\text{distinct limits}}.
\]
The theorem asserts this configuration is impossible in $\mathbb{R}$.
The failure mode is not merely exotic: it occurs in non-Hausdorff
topological spaces where distinct points cannot be separated by disjoint
open sets. The proof that it cannot occur in $\mathbb{R}$ is precisely
an argument that $\mathbb{R}$ has enough room to separate any two
distinct points by disjoint $\varepsilon$-neighborhoods.
\end{remark*}

\begin{remark*}[Interpretation]
The theorem licenses the notation $\displaystyle\lim_{n \to \infty} x_n
= L$: the definite article is justified precisely because the limit,
when it exists, is unique. Without uniqueness the limit would be a
relation rather than a function, and the notation would be ambiguous.
The proof strategy is a separation argument: if $L \neq K$ then
$|L - K| > 0$, and choosing $\varepsilon = |L - K|/2$ produces two
disjoint neighborhoods that the tail of $(x_n)$ cannot simultaneously
occupy. The Hausdorff remark in the failure mode block is the forward
pointer: when sequences are studied in general topological spaces,
uniqueness of limits is equivalent to the Hausdorff separation axiom.
\end{remark*}

\begin{remark*}[Dependencies]
\begin{itemize}
  \item \hyperref[def:convergent-sequence]{Convergence of a Sequence}
  \item \hyperref[thm:triangle-inequality]{Triangle inequality}
\end{itemize}
\end{remark*}





